\documentclass[a4paper,12pt]{article}
\usepackage[T2A]{fontenc}
\usepackage[utf8]{inputenc}
\usepackage[russian]{babel}
\usepackage{graphicx}
\usepackage{listingsutf8}   % UTF-8 в lstlisting (вместо listings)
\usepackage{xcolor}
\usepackage{geometry}
\usepackage{hyperref}
\usepackage{enumitem}
\usepackage{booktabs}
\usepackage{array}
\usepackage{amsmath}

\geometry{margin=2.5cm}

\hypersetup{
    colorlinks=true,
    linkcolor=blue,
    urlcolor=cyan
}

\lstset{
    language=C++,
    inputencoding=utf8,
    extendedchars=true,
    basicstyle=\ttfamily\small,
    keywordstyle=\color{blue}\bfseries,
    commentstyle=\color{gray}\itshape,
    stringstyle=\color{red!70!black},
    showstringspaces=false,
    breaklines=true,
    frame=single,
    numbers=left,
    numberstyle=\tiny\color{gray},
    backgroundcolor=\color{gray!8},
    tabsize=4,
    captionpos=b,
    xleftmargin=1.5em,
    framexleftmargin=1.5em
}

\title{\textbf{Лекции по C++}\\\large 2 семестр}
\author{Дмитрий Величко}
\date{Март 2026}

\begin{document}
\maketitle
\tableofcontents
\newpage

% -------------------------------------------------------
\section{Ключевое слово \texttt{throw}}

Рассмотрим функцию деления:

\begin{lstlisting}
#include <stdexcept>

int divide(int a, int b) {
    if (b == 0) {
        throw std::invalid_argument("division by zero");
    }
    return a / b;
}
\end{lstlisting}

При делении на ноль поведение программы неопределено, поэтому используется
ключевое слово \texttt{throw}. Бросать можно \textbf{любое выражение},
однако принято бросать объекты стандартных исключений из
\texttt{<stdexcept>}.

\texttt{throw} --- один из способов выйти из функции наряду с
\texttt{return}.

% -------------------------------------------------------
\section{Раскрутка стека}

При выбросе исключения происходит \textbf{раскрутка стека} (stack
unwinding):

\begin{enumerate}
    \item Уничтожаются локальные объекты текущей функции
          (в обратном порядке), вызываются их деструкторы.
    \item Управление передаётся вызывающей функции.
    \item Процесс повторяется вверх по стеку, пока не будет найден
          подходящий блок \texttt{catch}.
\end{enumerate}

Пример: если \texttt{main} вызвал \texttt{f1}, та вызвала
\texttt{f2}, \ldots, \texttt{f10}, и из \texttt{f10} брошено
исключение, то локальные объекты уничтожаются начиная с \texttt{f10}
и заканчивая \texttt{main}.

% -------------------------------------------------------
\section{try / catch}

\begin{lstlisting}
#include <iostream>
#include <stdexcept>

int main() {
    try {
        divide(1, 0);
    } catch (const std::exception& e) {
        std::cout << e.what() << "\n";
    }
}
\end{lstlisting}

\begin{itemize}
    \item Метод \texttt{what()} определён в \texttt{std::exception}
          и возвращает строку с описанием ошибки --- ту самую, что
          была передана при создании объекта исключения.
    \item Исключение следует ловить \textbf{по константной ссылке},
          чтобы избежать лишнего копирования и сохранить полиморфизм.
\end{itemize}

% -------------------------------------------------------
\section{Где встречаются исключения в стандартной библиотеке}

\begin{itemize}
    \item \texttt{vector::at()} --- бросает \texttt{std::out\_of\_range}
          при выходе за пределы (в отличие от \texttt{operator[]},
          который имеет неопределённое поведение).
    \item \texttt{operator new} --- бросает \texttt{std::bad\_alloc}
          при нехватке памяти.
    \item \texttt{dynamic\_cast<T\&>} --- бросает \texttt{std::bad\_cast}
          при неудачном приведении.
    \item \texttt{std::terminate()} --- вызывается при не пойманном
          исключении.
\end{itemize}

% -------------------------------------------------------
\section{Runtime-ошибки vs Исключения C++}

Это \textbf{принципиально разные вещи}.

\subsection{Три вида runtime-ошибок на уровне ОС}

\begin{center}
\begin{tabular}{lp{9cm}}
\toprule
\textbf{Сигнал} & \textbf{Типичные причины} \\
\midrule
\texttt{SIGSEGV} (Segfault)
    & Выход за пределы массива, разыменование \texttt{nullptr},
      переполнение стека \\
\addlinespace
\texttt{SIGFPE} (Floating Point Exception)
    & Деление целого на ноль, переполнение целых \\
\addlinespace
\texttt{SIGABRT} (Aborted)
    & Вызов \texttt{std::abort()}, \texttt{std::terminate()},
      падение \texttt{assert} \\
\bottomrule
\end{tabular}
\end{center}

\medskip
Эти ошибки \textbf{не являются} исключениями C++ и \textbf{не ловятся}
через \texttt{catch}:

\begin{lstlisting}
int arr[5];
try {
    arr[10] = 42;  // SIGSEGV -- catch не поймает
} catch (...) {
    // сюда мы НЕ попадём
}

try {
    int x = 1 / 0;  // SIGFPE -- тоже не исключение C++
} catch (...) {
    // и сюда не попадём
}
\end{lstlisting}

\textbf{Важно:} деление на ноль --- это исключение на уровне
\textit{процессора}, не языка. Процессор обращается к ОС, которая
посылает \texttt{SIGFPE}. Это совпадение терминологии, но разные
уровни абстракции.

\subsection{Пути к \texttt{std::terminate()}}

\begin{itemize}
    \item Не пойманное исключение (uncaught exception)
    \item Явный вызов \texttt{std::terminate()} или \texttt{std::abort()}
    \item Упавший \texttt{assert} (он напрямую вызывает \texttt{abort()})
    \item Проблема ромбовидного наследования при виртуальных деструкторах
\end{itemize}

С точки зрения ОС программа, завершившаяся из-за не пойманного C++-исключения,
получает вердикт \texttt{SIGABRT}.

Также стоит отметить, что \texttt{std::runtime\_error} ---
это \textit{частный случай} \texttt{std::exception}, то есть обычное
C++-исключение, не имеющее отношения к SIGFPE/SIGSEGV/SIGABRT.

% -------------------------------------------------------
\section{Копирование при броске исключения}

\begin{lstlisting}
#include <iostream>

struct A {
    A()         { std::cout << "A()\n"; }
    ~A()        { std::cout << "~A()\n"; }
    A(const A&) { std::cout << "A(copy)\n"; }
};

void f(int x) {
    A a;
    if (x == 0) {
        throw a;  // объект копируется в динамическую память
    }
}

int main() {
    try {
        f(0);
    } catch (...) {
        std::cout << "caught!\n";
    }
}
\end{lstlisting}

\textbf{Порядок вывода:}
\begin{enumerate}
    \item \texttt{A()} --- создание \texttt{a} на стеке
    \item \texttt{A(copy)} --- копирование в динамическую память при \texttt{throw}
    \item \texttt{\textasciitilde A()} --- уничтожение \texttt{a} при раскрутке стека
    \item \texttt{caught!}
    \item \texttt{\textasciitilde A()} --- уничтожение копии из динамической памяти
\end{enumerate}

\medskip
\begin{itemize}
    \item \texttt{catch(A a)} --- ещё одно копирование из дин.\ памяти на стек.
    \item \texttt{catch(A\& a)} --- работаем прямо в дин.\ памяти, без копирования.
\end{itemize}

\subsection{Где хранится объект исключения?}

Стандарт не специфицирует точное место хранения. На практике ---
в \textbf{динамической памяти}. Исключение: \texttt{std::bad\_alloc} ---
при нехватке памяти нельзя выделить новую, поэтому компилятор заранее
резервирует буфер в \textbf{статической памяти}
(\textit{emergency buffer}).

Исключения дороги по стоимости: выделение динамической памяти происходит
в любом случае. Поэтому их не стоит использовать в горячих путях кода.

\subsection{Повторный проброс}

\begin{lstlisting}
try {
    f(0);
} catch (const A& a) {
    throw;      // прокидываем то же исключение дальше (без копирования)
    // throw a; // создаём НОВЫЙ объект, возможно в другом месте памяти
}
\end{lstlisting}

Если исключение не поймано нигде --- гарантируется только вызов
\texttt{std::terminate()}. Деструкторы при этом могут не вызваться.

% -------------------------------------------------------
\section{Приведение типов в \texttt{catch}}

В \texttt{catch} \textbf{не работают} обычные неявные преобразования
и правила перегрузки. Доступны только:
\begin{itemize}
    \item Добавление \texttt{const} / приведение к ссылке
    \item Приведение типа потомка к типу родителя
\end{itemize}

\begin{lstlisting}
try {
    throw 1;  // int
} catch (double d) {
    std::cout << "double\n";    // не поймает: int -> double не работает
} catch (long long l) {
    std::cout << "long long\n"; // не поймает: int -> long long не работает
} catch (...) {
    std::cout << "other\n";     // поймает
}
\end{lstlisting}

\subsection{Наследование в \texttt{catch}}

Блоки \texttt{catch} проверяются \textbf{сверху вниз}, побеждает
первый подходящий:

\begin{lstlisting}
struct Mom {};
struct Son : Mom {};

int main() {
    try {
        throw Son{};
    } catch (const Mom& m) {
        std::cout << "Mom\n"; // поймает: Son приводится к Mom
    } catch (const Son& s) {
        std::cout << "Son\n"; // сюда не дойдём
    } catch (...) {
        std::cout << "other\n";
    }
}
\end{lstlisting}

\textbf{Вывод:} \texttt{Mom}. Именно поэтому обработчики потомков
нужно ставить \textbf{перед} обработчиками родителей, иначе они
никогда не сработают.

% -------------------------------------------------------
\newpage
\section{Контейнеры и итераторы}

\subsection{\texttt{std::vector}: реализация изнутри}

Начнём с устройства \texttt{std::vector} и в первую очередь рассмотрим
реализацию метода \texttt{push\_back}.

Структура вектора состоит из трёх полей: указателя на данные \texttt{arr\_},
размера \texttt{sz\_} и ёмкости \texttt{cap\_}.

\subsubsection{Почему нельзя писать \texttt{new T[newcap]}}

При реаллокации наивная реализация могла бы выглядеть так:

\begin{lstlisting}
T* newarr = new T[newcap]; // проблема!
\end{lstlisting}

Это не работает по двум причинам:
\begin{itemize}
    \item Мы не знаем, есть ли у \texttt{T} конструктор по умолчанию.
    \item Нам нужно выделить \textit{сырую} память и заполнять её по мере
          надобности, а не создавать все объекты сразу.
\end{itemize}

\subsubsection{Placement new}

Решение --- выделить сырую память через \texttt{char}, а объекты
создавать по месту с помощью \textbf{placement new}:

\begin{lstlisting}
new (ptr) T(args...);  // создаёт объект T по адресу ptr
\end{lstlisting}

Почему нельзя просто написать \texttt{newarr[i] = arr\_[i]}?
Потому что \texttt{newarr} указывает на сырые байты --- объекта
\texttt{T} там ещё нет, и присваивать можно только существующему объекту.
\texttt{memcpy} тоже не подходит: он не работает корректно, если объект
содержит указатель или ссылку на своё собственное поле
(например, \texttt{std::string} с small-string optimization или
декартово дерево).

\subsubsection{Почему нельзя \texttt{delete[] arr\_}}

При освобождении старой памяти нельзя писать \texttt{delete[] arr\_}:
это вызовет деструктор для каждого из \texttt{cap\_} слотов, хотя
реальных объектов там только \texttt{sz\_}. Нужно:
\begin{enumerate}
    \item вручную вызвать деструктор для каждого существующего объекта;
    \item освободить сырую память через
          \texttt{delete[] reinterpret\_cast<char*>(arr\_)}.
\end{enumerate}

\subsubsection{Реализация \texttt{reserve}}

\begin{lstlisting}
template<typename T>
class vector {
    T*     arr_;
    size_t sz_;
    size_t cap_;

    void reserve(size_t newcap) {
        T* newarr = reinterpret_cast<T*>(new char[newcap * sizeof(T)]);
        size_t i = 0;
        try {
            for (; i < sz_; ++i)
                new (newarr + i) T(arr_[i]); // placement new
        } catch (...) {
            for (size_t k = 0; k < i; ++k)
                (newarr + k)->~T();
            delete[] reinterpret_cast<char*>(newarr);
            throw;
        }
        for (size_t k = 0; k < sz_; ++k)
            (arr_ + k)->~T();
        delete[] reinterpret_cast<char*>(arr_);
        arr_ = newarr;
        cap_ = newcap;
    }
\end{lstlisting}

\textbf{Обработка исключений в \texttt{reserve}:}
\begin{itemize}
    \item Если \texttt{new char[...]} бросает \texttt{std::bad\_alloc}
          --- мы просто не выделили память, ничего не испорчено, вылетаем наверх.
    \item Если placement new бросает исключение в конструкторе на шаге \texttt{i}
          --- нужно уничтожить уже созданные объекты (с 0 по \texttt{i-1}),
          освободить \texttt{newarr} и перебросить исключение.
\end{itemize}

\subsubsection{Реализация \texttt{push\_back}}

\begin{lstlisting}
    void push_back(const T& value) {
        if (sz_ == cap_)
            reserve(cap_ > 0 ? cap_ * 2 : 1);
        new (arr_ + sz_) T(value);
        ++sz_;
    }
}; // конец class vector
\end{lstlisting}

% -------------------------------------------------------
\subsection{\texttt{vector<bool>}: специализация с битовым хранением}

\texttt{vector<bool>} --- это печально известная специализация,
которая ломает многие интуитивные ожидания от вектора.

\textbf{Идея:} хранить элементы побитово, упаковывая по 8 штук в \texttt{char}.
Это экономит память, но делает невозможным взять ссылку \texttt{bool\&} на
отдельный элемент --- ссылка не может указывать на отдельный бит.

\subsubsection{Прокси-класс \texttt{BitReference}}

Чтобы сохранить синтаксис \texttt{v[i] = true}, вводится прокси-класс:

\begin{lstlisting}
template<>
class vector<bool> {
    char*  arr_;
    size_t sz_;
    size_t cap_;

    struct BitReference {
        char*   cell;
        uint8_t idx;

        BitReference(char* cell, uint8_t idx) : cell(cell), idx(idx) {}

        void operator=(bool b) {
            if (b) *cell |=  (1 << idx);
            else   *cell &= ~(1 << idx);
        }

        operator bool() const {
            return *cell & (1 << idx);
        }
    };

public:
    BitReference operator[](size_t idx) {
        return BitReference(arr_ + idx / 8, idx % 8);
    }
};
\end{lstlisting}

\subsubsection{Трюк Ильи Мещерина: узнать тип выражения}

\begin{lstlisting}
template<typename T>
class Debug {
    Debug(T) = delete;
};

int main() {
    std::vector<bool> v(10);
    v[5] = true;
    // Debug d(v[5]); -- компилятор скажет тип: std::__bit_reference
}
\end{lstlisting}

\textbf{Почему это проблема:} \texttt{vector<bool>} нарушает концепцию
контейнера --- из него нельзя взять настоящую ссылку на элемент
(\texttt{bool\& r = v[5]} не скомпилируется), и многие шаблонные
алгоритмы ломаются при работе с ним. Поэтому для хранения битовых флагов
часто предпочитают \texttt{std::vector<char>} или \texttt{std::bitset}.

% -------------------------------------------------------
\newpage
\section{Библиотека \texttt{<algorithm>}}

\subsection{Зачем она вообще нужна}

Стандартная библиотека предоставляет \textbf{готовые алгоритмы},
которые работают с любыми контейнерами через \textbf{итераторы}.
Ключевая идея: алгоритм не знает, что за контейнер перед ним ---
он получает пару итераторов \texttt{[begin, end)} и работает с ними.

\begin{lstlisting}
#include <algorithm>
#include <vector>

int main() {
    int arr[] = {5, 3, 1, 4, 2};
    std::sort(arr, arr + 5);

    std::vector<int> v = {5, 3, 1, 4, 2};
    std::sort(v.begin(), v.end());
}
\end{lstlisting}

% -------------------------------------------------------
\subsection{\texttt{std::sort}}

Сортирует диапазон \textbf{на месте}. Сложность: $O(n \log n)$ (introsort).

\begin{lstlisting}
std::vector<int> v = {5, 3, 1, 4, 2};
std::sort(v.begin(), v.end()); // по возрастанию

std::sort(v.begin(), v.end(), [](int a, int b) {
    return a > b; // по убыванию
});
\end{lstlisting}

\textbf{Важно:} компаратор должен задавать \textit{строгий слабый порядок}
(\texttt{<}, а не \texttt{<=}). Если написать \texttt{a >= b} --- UB.

% -------------------------------------------------------
\subsection{\texttt{std::find} и \texttt{std::find\_if}}

Линейный поиск. Возвращает итератор на первый подходящий элемент,
или \texttt{last} если ничего не нашли.

\begin{lstlisting}
std::vector<int> v = {1, 2, 3, 4, 5};

auto it = std::find(v.begin(), v.end(), 3);
if (it != v.end()) std::cout << *it; // 3

auto it2 = std::find_if(v.begin(), v.end(), [](int x) {
    return x % 2 == 0; // первый чётный -> 2
});
\end{lstlisting}

\texttt{end()} как «не найдено» --- универсальное соглашение, потому что
алгоритм не знает о типе контейнера и не может вернуть \texttt{-1}.

% -------------------------------------------------------
\subsection{\texttt{std::count} и \texttt{std::count\_if}}

\begin{lstlisting}
std::vector<int> v = {1, 2, 2, 3, 2, 4};
int cnt  = std::count(v.begin(), v.end(), 2);                      // 3
int even = std::count_if(v.begin(), v.end(), [](int x){ return x%2==0; }); // 4
\end{lstlisting}

% -------------------------------------------------------
\subsection{\texttt{std::for\_each} и \texttt{std::transform}}

\begin{lstlisting}
std::vector<int> v = {1, 2, 3, 4, 5};

std::for_each(v.begin(), v.end(), [](int& x) { x *= 2; });
// v == {2, 4, 6, 8, 10}

std::vector<int> out(v.size());
std::transform(v.begin(), v.end(), out.begin(), [](int x) {
    return x * x;
});
// out == {4, 16, 36, 64, 100}
\end{lstlisting}

% -------------------------------------------------------
\subsection{\texttt{std::copy} и \texttt{std::back\_inserter}}

\begin{lstlisting}
std::vector<int> src = {1, 2, 3, 4, 5};
std::vector<int> evens;

std::copy_if(src.begin(), src.end(), std::back_inserter(evens),
             [](int x) { return x % 2 == 0; });
// evens == {2, 4}
\end{lstlisting}

\texttt{std::back\_inserter} --- итератор-адаптер, который при присваивании
вызывает \texttt{push\_back}. Нужен, когда выходной контейнер изначально пуст.

% -------------------------------------------------------
\subsection{Идиома erase-remove}

\texttt{std::remove} \textbf{не удаляет} элементы --- он лишь сдвигает
нужные в начало и возвращает итератор на новый логический конец.
Реальное удаление делает \texttt{erase}:

\begin{lstlisting}
std::vector<int> v = {1, 2, 3, 2, 4, 2, 5};
v.erase(std::remove(v.begin(), v.end(), 2), v.end());
// v == {1, 3, 4, 5}
\end{lstlisting}

% -------------------------------------------------------
\subsection{\texttt{std::accumulate} из \texttt{<numeric>}}

\begin{lstlisting}
#include <numeric>
std::vector<int> v = {1, 2, 3, 4, 5};
int sum     = std::accumulate(v.begin(), v.end(), 0);             // 15
int product = std::accumulate(v.begin(), v.end(), 1,
              [](int acc, int x) { return acc * x; });            // 120
\end{lstlisting}

% -------------------------------------------------------
\newpage
\section{Дек (\texttt{std::deque})}

\subsection{Инвалидация указателей и ссылок в \texttt{vector}}

\begin{lstlisting}
std::vector<int> v(10);
int* p = &v[5];
int& x = v[5];
v.push_back(1);      // реаллокация!
std::cout << *p;     // UB -- pointer invalidation
std::cout << x;      // UB -- reference invalidation
\end{lstlisting}

При \texttt{push\_back} вектор мог реаллоцироваться. Старый кусок
памяти нам больше не принадлежит --- все указатели и ссылки на элементы,
полученные до реаллокации, недействительны.

\subsection{Два главных отличия дека от вектора}

\begin{enumerate}
    \item \texttt{std::deque} \textbf{не инвалидирует} указатели и ссылки
          на существующие элементы при добавлении новых.
    \item Дек умеет \texttt{push\_front()} и \texttt{pop\_front()} за $O(1)$.
\end{enumerate}

\subsection{Внутреннее устройство дека}

Дек хранит \textbf{массив указателей на бакеты} фиксированного размера ---
\texttt{T** arr\_}.

\begin{itemize}
    \item При \texttt{push\_back}: идём вправо по бакету; когда бакет
          заполнен --- выделяем новый бакет справа.
    \item При \texttt{push\_front}: идём влево; когда бакет заполнен ---
          выделяем новый бакет слева и кладём элемент с его конца.
    \item Дек хранит две пары: (номер бакета начала, индекс внутри него)
          и (номер бакета конца, индекс внутри него).
\end{itemize}

\textbf{Почему не инвалидируются указатели?} При переполнении внешнего
массива реаллоцируется только он --- перепривязываются указатели на бакеты.
Сами бакеты с элементами остаются на месте.

\textbf{\texttt{operator[]} за $O(1)$:} зная индекс начала и размер бакета,
пересчитываем номер бакета и смещение за константное время.

\medskip
У дека нет \texttt{reserve()} и \texttt{shrink\_to\_fit()} --- элементы
разбросаны по бакетам, и их нельзя реаллоцировать одним куском.

% -------------------------------------------------------
\section{Адаптеры контейнеров}

\texttt{std::stack}, \texttt{std::queue} и \texttt{std::priority\_queue}
--- это \textbf{не контейнеры}, а \textbf{адаптеры}: они оборачивают
контейнер и предоставляют ограниченный интерфейс поверх него.

\begin{lstlisting}
std::stack<int>                    s;   // по умолчанию на deque
std::stack<int, std::vector<int>>  s2;  // на vector
std::stack<int, std::list<int>>    s3;  // на list
\end{lstlisting}

\subsection{\texttt{std::stack}}

\begin{lstlisting}
template<typename T, typename Container = std::deque<T>>
class stack {
    Container c_;
public:
    void   push(const T& val) { c_.push_back(val); }
    void   pop()              { c_.pop_back(); }
    T&     top()              { return c_.back(); }
    bool   empty() const      { return c_.empty(); }
    size_t size()  const      { return c_.size(); }
};
\end{lstlisting}

\textbf{Почему \texttt{pop()} возвращает \texttt{void}?}
Если бы \texttt{pop()} возвращал значение, пришлось бы скопировать
элемент перед удалением. Если конструктор копирования бросит исключение
--- элемент уже удалён и значение потеряно. Разделение на \texttt{top()}
(посмотреть) и \texttt{pop()} (удалить) даёт exception safety.
Вернуть ссылку тоже нельзя: после \texttt{pop()} она стала бы битой.

% -------------------------------------------------------
\newpage
\section{Итераторы и их категории}

\subsection{Что такое итератор}

Итератор --- это тип, который позволяет \textbf{обходить последовательность}.
Это обобщение указателя: его можно разыменовать (\texttt{*it}),
инкрементировать (\texttt{++it}) и сравнивать (\texttt{it != end}).
Обычный сырой указатель --- полноценный итератор.

Range-based for --- синтаксический сахар:
\begin{lstlisting}
for (int x : s) { ... }
// разворачивается в:
for (auto it = s.begin(); it != s.end(); ++it) { ... }
\end{lstlisting}
Достаточно реализовать \texttt{begin()} и \texttt{end()}.

\subsection{Категории итераторов}

\begin{center}
\begin{tabular}{lp{5.2cm}p{4cm}}
\toprule
\textbf{Категория} & \textbf{Что умеет} & \textbf{Примеры} \\
\midrule
\texttt{InputIterator}
    & \texttt{*it}, \texttt{++it}, \texttt{==}, \texttt{!=}
    & потоковые итераторы \\
\addlinespace
\texttt{ForwardIterator}
    & + гарантия повторного прохода
    & \texttt{unordered\_set}, \texttt{forward\_list} \\
\addlinespace
\texttt{BidirectionalIterator}
    & + \texttt{-{}-it}
    & \texttt{list}, \texttt{set}, \texttt{map} \\
\addlinespace
\texttt{RandomAccessIterator}
    & + \texttt{it+=n}, \texttt{it[n]}, \texttt{<}, \texttt{>}
    & \texttt{deque} \\
\addlinespace
\texttt{ContiguousIterator}
    & + гарантия непрерывности памяти
    & \texttt{vector}, \texttt{array}, указатель \\
\bottomrule
\end{tabular}
\end{center}

\medskip
Разница между \texttt{Input} и \texttt{Forward}: \texttt{InputIterator}
не гарантирует, что при повторном проходе увидим те же данные (поток ---
прочитанное не вернёшь). \texttt{ForwardIterator} эту гарантию даёт.

\subsection{\texttt{iterator\_traits}}

Для получения информации об итераторе внутри шаблонной функции:

\begin{lstlisting}
template<typename It>
void foo(It begin, It end) {
    typename std::iterator_traits<It>::value_type x = *begin;
    using Cat = typename std::iterator_traits<It>::iterator_category;
}
\end{lstlisting}

Доступные поля: \texttt{value\_type}, \texttt{reference}, \texttt{pointer},
\texttt{difference\_type}, \texttt{iterator\_category}.

Почему не \texttt{auto x = *begin}? Для \texttt{vector<bool>} \texttt{auto}
выведет \texttt{BitReference}, а не \texttt{bool}. \texttt{value\_type}
вернёт правильный тип.

\subsection{Реализация \texttt{std::distance}}

\begin{lstlisting}
template<typename It>
typename std::iterator_traits<It>::difference_type
distance(It first, It last) {
    if constexpr (std::is_same_v<
        typename std::iterator_traits<It>::iterator_category,
        std::random_access_iterator_tag>)
    {
        return last - first;  // O(1)
    }
    typename std::iterator_traits<It>::difference_type i = 0;
    for (; first != last; ++first, ++i) {}
    return i;                 // O(n)
}
\end{lstlisting}

\texttt{if constexpr} вычисляется на этапе компиляции: невыбранная ветка
не компилируется вообще. Это позволяет писать \texttt{last - first} не
боясь, что для non-random-access итератора это не скомпилируется.

\subsection{Реализация итератора вектора}

\begin{lstlisting}
template<typename T>
class Vector {
public:
    class Iterator {
    public:
        using iterator_category = std::random_access_iterator_tag;
        using value_type        = T;
        using difference_type   = std::ptrdiff_t;
        using pointer           = T*;
        using reference         = T&;

        explicit Iterator(pointer ptr) : ptr_(ptr) {}

        reference operator*()  const { return *ptr_; }
        pointer   operator->() const { return  ptr_; }

        Iterator& operator++()    { ++ptr_; return *this; }
        Iterator  operator++(int) { Iterator tmp = *this; ++ptr_; return tmp; }
        Iterator& operator--()    { --ptr_; return *this; }
        Iterator  operator--(int) { Iterator tmp = *this; --ptr_; return tmp; }

        Iterator  operator+(difference_type n) const { return Iterator(ptr_ + n); }
        Iterator  operator-(difference_type n) const { return Iterator(ptr_ - n); }
        difference_type operator-(const Iterator& o) const { return ptr_ - o.ptr_; }
        Iterator& operator+=(difference_type n) { ptr_ += n; return *this; }
        Iterator& operator-=(difference_type n) { ptr_ -= n; return *this; }
        reference operator[](difference_type n) const { return ptr_[n]; }

        bool operator==(const Iterator& o) const { return ptr_ == o.ptr_; }
        bool operator!=(const Iterator& o) const { return ptr_ != o.ptr_; }
        bool operator< (const Iterator& o) const { return ptr_ <  o.ptr_; }
        bool operator> (const Iterator& o) const { return ptr_ >  o.ptr_; }
        bool operator<=(const Iterator& o) const { return ptr_ <= o.ptr_; }
        bool operator>=(const Iterator& o) const { return ptr_ >= o.ptr_; }

    private:
        pointer ptr_;
    };

    Vector() : data_(nullptr), size_(0), capacity_(0) {}
    ~Vector() { delete[] data_; }

    void push_back(const T& val) {
        if (size_ == capacity_) reserve(capacity_ == 0 ? 1 : capacity_ * 2);
        data_[size_++] = val;
    }

    T&          operator[](std::size_t i) { return data_[i]; }
    std::size_t size()     const { return size_; }
    std::size_t capacity() const { return capacity_; }

    Iterator begin() { return Iterator(data_); }
    Iterator end()   { return Iterator(data_ + size_); }

private:
    void reserve(std::size_t new_cap) {
        T* tmp = new T[new_cap];
        for (std::size_t i = 0; i < size_; ++i) tmp[i] = data_[i];
        delete[] data_;
        data_ = tmp; capacity_ = new_cap;
    }

    T* data_; std::size_t size_; std::size_t capacity_;
};
\end{lstlisting}

Пять \texttt{type alias}-ов в \texttt{Iterator} --- контракт с
\texttt{iterator\_traits}. Без них \texttt{std::sort} и
\texttt{std::distance} не смогут узнать категорию итератора.

% -------------------------------------------------------
\section{Потоковые итераторы}

\subsection{\texttt{std::istream\_iterator}}

\begin{lstlisting}
#include <iterator>
#include <vector>

int main() {
    std::vector<int> v(
        std::istream_iterator<int>(std::cin),
        std::istream_iterator<int>()   // итератор-конец (EOF)
    );
}
\end{lstlisting}

Это \textbf{InputIterator} --- одноразовый: прочитанное назад не вернуть.
\texttt{istream\_iterator<int>()} без аргументов --- итератор конца потока.

\begin{lstlisting}
std::istream_iterator<int> in(std::cin), eof;
int sum = std::accumulate(in, eof, 0); // сумма всех чисел из stdin
\end{lstlisting}

\subsection{\texttt{std::ostream\_iterator}}

\begin{lstlisting}
std::vector<int> v = {1, 2, 3, 4, 5};

std::copy(v.begin(), v.end(),
          std::ostream_iterator<int>(std::cout, " "));
// вывод: 1 2 3 4 5

std::transform(v.begin(), v.end(),
               std::ostream_iterator<int>(std::cout, "\n"),
               [](int x) { return x * x; });
\end{lstlisting}

\texttt{ostream\_iterator} --- \textbf{OutputIterator}: присваивание
\texttt{*it = val} вызывает \texttt{os << val << delim}.

\subsection{Итераторы файловых потоков}

\begin{lstlisting}
#include <fstream>
#include <iterator>
#include <algorithm>

int main() {
    std::ifstream fin("input.txt");
    std::ofstream fout("output.txt");

    std::vector<int> v(std::istream_iterator<int>(fin),
                       std::istream_iterator<int>());
    std::sort(v.begin(), v.end());
    std::copy(v.begin(), v.end(),
              std::ostream_iterator<int>(fout, "\n"));
}
\end{lstlisting}

Алгоритмы не знают, откуда данные читаются и куда пишутся --- они
работают с абстракцией итератора. В этом и есть смысл всей иерархии.

% -------------------------------------------------------
\newpage
\section{Аллокаторы и управление памятью}

\subsection{Зачем аллокатор нужен контейнеру}

У каждого контейнера есть шаблонный параметр аллокатора:

\begin{lstlisting}
std::list<int>    // = std::list<int, std::allocator<int>>
std::vector<int>  // = std::vector<int, std::allocator<int>>
\end{lstlisting}

Аллокатор --- это \textbf{промежуточный класс между контейнером и
оператором \texttt{new}}. Стек вызовов:

\begin{center}
\texttt{Container} $\to$ \texttt{Allocator} $\to$
\texttt{operator new} $\to$ \texttt{malloc} $\to$ \texttt{OS}
\end{center}

\subsection{Устройство стандартного аллокатора}

У аллокатора четыре метода:

\begin{lstlisting}
template<typename T>
struct allocator {
    using value_type = T;

    T* allocate(std::size_t n) {
        if (n == 0) return nullptr;
        return static_cast<T*>(::operator new(n * sizeof(T)));
    }

    void deallocate(T* p, std::size_t) {
        ::operator delete(p);
    }

    template<typename... Args>
    void construct(T* p, Args&&... args) {
        ::new (static_cast<void*>(p)) T(std::forward<Args>(args)...);
    }

    void destroy(T* p) {
        p->~T();
    }
};
\end{lstlisting}

\begin{itemize}
    \item \texttt{allocate} --- выделяет сырую память, конструкторы не вызываются.
    \item \texttt{deallocate} --- освобождает память, деструкторы не вызываются.
    \item \texttt{construct} --- вызывает конструктор по адресу (placement new).
    \item \texttt{destroy} --- явно вызывает деструктор, память не освобождает.
\end{itemize}

\subsection{Проблема: \texttt{list} аллоцирует не \texttt{T}}

\texttt{std::list<int>} хранит не просто \texttt{int}, а узлы:

\begin{lstlisting}
template<typename T>
struct ListNode { T value; ListNode* prev; ListNode* next; };
\end{lstlisting}

Но аллокатор шаблонизирован по \texttt{T = int}. Решение комитета ---
\textbf{rebind}: вложенный шаблон в аллокаторе, который перепривязывает
его к другому типу:

\begin{lstlisting}
template<typename T>
struct allocator {
    template<typename U>
    struct rebind { using other = allocator<U>; };
};

// list внутри делает:
using NodeAlloc = typename Alloc::template rebind<ListNode<T>>::other;
NodeAlloc node_alloc_;  // аллокатор для узлов, а не для T
\end{lstlisting}

\subsection{\texttt{allocator\_traits}}

\texttt{std::allocator\_traits<Alloc>} --- адаптер с единым интерфейсом
к любому аллокатору и разумными дефолтами для отсутствующих методов.

\begin{lstlisting}
template<typename Alloc>
struct allocator_traits {
    using value_type      = typename Alloc::value_type;
    using pointer         = value_type*;
    using size_type       = std::size_t;
    using difference_type = std::ptrdiff_t;

    template<typename U>
    using rebind_alloc = typename Alloc::template rebind<U>::other;

    static pointer allocate(Alloc& a, size_type n)
        { return a.allocate(n); }
    static void deallocate(Alloc& a, pointer p, size_type n)
        { a.deallocate(p, n); }

    // если у Alloc нет construct -- placement new по умолчанию
    template<typename T, typename... Args>
    static void construct(Alloc&, T* p, Args&&... args)
        { ::new (static_cast<void*>(p)) T(std::forward<Args>(args)...); }

    // если у Alloc нет destroy -- явный деструктор по умолчанию
    template<typename T>
    static void destroy(Alloc&, T* p) { p->~T(); }

    // политики распространения аллокатора
    using propagate_on_container_copy_assignment = std::false_type;
    using propagate_on_container_move_assignment = std::false_type;
    using propagate_on_container_swap            = std::false_type;
};
\end{lstlisting}

\textbf{Главная идея}: контейнер должен работать с аллокатором
\textit{только через \texttt{allocator\_traits}}. Тогда кастомный
аллокатор может не реализовывать \texttt{construct}/\texttt{destroy}
и всё равно работать.

\subsubsection{Политики распространения аллокатора}

\begin{center}
\begin{tabular}{lp{7.5cm}}
\toprule
\textbf{Флаг} & \textbf{Что означает} \\
\midrule
\texttt{propagate\_on\_container\_copy\_assignment}
    & Копировать аллокатор при \texttt{operator=} копирования? \\
\addlinespace
\texttt{propagate\_on\_container\_move\_assignment}
    & Копировать аллокатор при \texttt{operator=} перемещения? \\
\addlinespace
\texttt{propagate\_on\_container\_swap}
    & Менять аллокаторы местами при \texttt{swap}? \\
\bottomrule
\end{tabular}
\end{center}

По умолчанию все три --- \texttt{false\_type}. Это важно для stateful
аллокаторов: если два вектора используют разные арены, при свопе нужно
либо перетащить элементы, либо поменять аллокаторы.

\subsection{Allocator-aware контейнер}

\begin{lstlisting}
template<typename T, typename Alloc = std::allocator<T>>
class vector {
    using AllocTraits = std::allocator_traits<Alloc>;

    Alloc  alloc_;
    T*     arr_;
    size_t sz_;
    size_t cap_;

public:
    explicit vector(const Alloc& a = Alloc())
        : alloc_(a), arr_(nullptr), sz_(0), cap_(0) {}

    void push_back(const T& val) {
        if (sz_ == cap_) reserve(cap_ ? cap_ * 2 : 1);
        AllocTraits::construct(alloc_, arr_ + sz_, val);
        ++sz_;
    }

    ~vector() {
        for (size_t i = 0; i < sz_; ++i)
            AllocTraits::destroy(alloc_, arr_ + i);
        AllocTraits::deallocate(alloc_, arr_, cap_);
    }

private:
    void reserve(size_t newcap) {
        T* newarr = AllocTraits::allocate(alloc_, newcap);
        size_t i = 0;
        try {
            for (; i < sz_; ++i)
                AllocTraits::construct(alloc_, newarr + i, arr_[i]);
        } catch (...) {
            for (size_t k = 0; k < i; ++k)
                AllocTraits::destroy(alloc_, newarr + k);
            AllocTraits::deallocate(alloc_, newarr, newcap);
            throw;
        }
        for (size_t k = 0; k < sz_; ++k)
            AllocTraits::destroy(alloc_, arr_ + k);
        AllocTraits::deallocate(alloc_, arr_, cap_);
        arr_ = newarr; cap_ = newcap;
    }
};
\end{lstlisting}

Stateless аллокатор (как \texttt{std::allocator}) не занимает места
благодаря \textbf{EBO} (Empty Base Optimization).

\subsection{Перегрузка \texttt{operator new} и \texttt{operator delete}}

\subsubsection{Глобальная перегрузка}

\begin{lstlisting}
void* operator new(std::size_t size) {
    void* ptr = std::malloc(size);
    if (!ptr) throw std::bad_alloc{};
    return ptr;
}
void operator delete(void* ptr) noexcept { std::free(ptr); }

// C++14: sized delete
void operator delete(void* ptr, std::size_t size) noexcept { std::free(ptr); }
\end{lstlisting}

Затрагивает весь процесс. Используется для профилировщиков памяти и
детекторов утечек.

\subsubsection{Классовая перегрузка}

\begin{lstlisting}
struct Node {
    int value; Node* next;

    void* operator new(std::size_t size) {
        return ::operator new(size); // делегируем глобальному
    }
    void operator delete(void* ptr) noexcept {
        ::operator delete(ptr);
    }
};
\end{lstlisting}

При \texttt{new Node} компилятор сначала ищет \texttt{operator new}
в классе, и только если не нашёл --- берёт глобальный.

\subsubsection{Когда использовать что}

\begin{center}
\begin{tabular}{lp{7.5cm}}
\toprule
\textbf{Инструмент} & \textbf{Когда применять} \\
\midrule
Кастомный аллокатор
    & Своя стратегия памяти для конкретного контейнера \\
\addlinespace
Классовая перегрузка \texttt{new/delete}
    & Свой пул для конкретного класса (например, узлы дерева) \\
\addlinespace
Глобальная перегрузка \texttt{new/delete}
    & Профилирование, детектирование утечек для всей программы \\
\bottomrule
\end{tabular}
\end{center}

\end{document}
